\section{Discussion}
\label{sec:discussion}

\subsection{What Binds Performance Now}

\paragraph{The front end, not the code.}  The clearest lesson of the
coding work is that FEC choice stopped mattering before code quality
did: LDPC's 0.7\,dB AWGN advantage compresses to parity end-to-end
because the demodulator's LLR quality is the binding constraint.  The
monotone recalibration recovered 1.5--2\,dB, but a standards-grade
LDPC matrix (802.11n/NR-style~\cite{ieee80211}) \emph{plus} fully
calibrated LLRs would be needed to realise the remaining
$\sim$2\,dB.  Any further sensitivity work should start at the LLR
front end, not the decoder.

\paragraph{Synchronization is the sensitivity frontier.}  Both the
original reproduction gap and the fade-loss analysis point the same
way: at the edge, frames are lost to detection and timing, not to
decoding.  HARQ helps only in the noise-limited regime for exactly
this reason --- a sync-level loss leaves no LLRs to combine.  The
two-board stand made the same point at the \emph{other} end of the
SNR range: broadcast, which cannot retransmit, lost one group in 39 on
a wire reporting $+24.7$\,dB --- a missed acquisition with 25\,dB of
margin, invisible to every acknowledged mode because ARQ had been
quietly repairing it (Section~\ref{sec:cport_bcast}).

\paragraph{Physics caps EXTREME.}  The $64\times$ mode's 0.69\,s
coherent integration assumes a channel coherent over that window
($\lesssim$0.1\,Hz Doppler spread).  This is a physical ceiling, not
an implementation one, and it is why the mode ladder --- rather than
ever-longer integration --- is the right architecture.

\paragraph{Feedback latency at the floor.}  Once a direction falls to
EXTREME, the adaptation loop period is bounded by 25--45\,s frame air
times.  The controller is designed around this (fade-aware filtering,
loss fallback, silence decay), but no tuning can make a 7.8\,bit/s
control loop fast.

\subsection{Limitations}

The channel model, while faithful to the article (and extended with
Rayleigh fading and a physical LO model), is still a simulation; the
measured sensitivities should be planned around with $\sim$3\,dB of
fading margin ($-14\ldots-15$\,dB for EXTREME).  The 16-QAM rungs run
without LLR recalibration (the map is BPSK-trained); a
modulation-specific map is straightforward but unmeasured.  The LDPC
matrix is a from-scratch IRA construction, not a standards matrix.
The project also has no over-the-air validation yet --- the WAV
recordings are transceiver-ready precisely to enable it.  The voice
evaluation of \S\ref{sec:cport_voice} is bounded the same way: the
codec's catastrophic-failure rate was measured on read speech, not on
shack audio, and one non-English trial is an anecdote rather than a
rate.

\subsection{Future Work}

In rough order of expected value: over-the-air testing with real SSB
rigs using the recorded WAVs; a standards-grade LDPC matrix plus
QAM-specific LLR recalibration; porting the fixed receiver's
slew-limited tracker back to the float model (it is measurably
better); pilot reduction or 2-D channel estimation to reclaim carrier
capacity; per-carrier bit-loading on the measured channel estimate;
frame aggregation in the ARQ to amortise preamble overhead at high
rungs; a host-side voice mode over the 250\,bit/s codec evaluated in
\S\ref{sec:cport_voice}, whose remaining question is a failure rate
rather than a bitrate; a 46.875\,Hz-spacing variant for very stable channels; and
shrinking EXTREME's frequency-search range after AFC netting
converges (compute reduction on embedded targets).
